% !TEX root = template.tex

\section{Introduction}
\label{sec:introduction}

Wireless sensing turns existing WiFi infrastructure into a device-free
perception system: a moving person perturbs the multipath propagation of the
radio signal, and those perturbations are recorded in the Channel State
Information (CSI) that commodity network cards expose. Human Activity
Recognition (HAR) from CSI needs no wearable, no camera and no cooperation
from the subject, which makes it attractive for assisted living, surveillance
and smart-home applications~\cite{Wang-2019}. The micro-Doppler signature ---
the distribution of reflected power over velocity and time --- is among the
most informative representations derived from CSI, and is our input throughout.

What separates CSI HAR from deployment is not in-distribution accuracy, which
is by now high, but \emph{generalization across environments}: a model trained
in one room entangles the activity with that room's static geometry. The
literature answers with two representation-shaping objectives ---
\emph{adversarial domain invariance}, where a domain classifier behind a
gradient reversal layer (GRL)~\cite{Ganin-2016} pushes the encoder to discard
environment-specific cues, and \emph{supervised contrastive} (SupCon)
learning~\cite{Khosla-2020}, which promises a more transferable embedding
geometry. Both are usually reported on the benchmarks where they helped, each
with its own backbone and protocol; whether they help on a given dataset, and
what must be true of that dataset for them to help at all, is assumed far more
often than it is tested.

We test exactly that, on the SHARP micro-Doppler
dataset~\cite{Meneghello-2023} under a leave-one-environment-out (LOEO)
protocol. Front-end, windowing, backbone, sampler, optimizer, seed and
evaluation harness are held identical, so the training objective is the only
variable; checkpoints are selected on validation alone, the test set is
evaluated exactly once from a frozen row list, and differences are assessed
with a paired bootstrap that treats the recording, not the window, as the
statistical unit. The answer is a controlled negative result, and its value
lies in the explanation: probes on the frozen encoders show that no domain
attribute is linearly readable from the features an adversary would consume,
so on data with this structure the adversarial route could not have worked.
That precondition costs one linear probe on cached features, and we argue it
should be measured before any invariance objective is trained.

\textbf{Contributions.} The key contributions of this report are:
\begin{itemize}
\item A controlled LOEO comparison of CE, CE+GRL and SupCon for
cross-environment CSI HAR under one apparatus and a single-use test protocol;
our CE residual network reaches macro-F1 $0.804$ on the unseen environment and
no variant improves on it.
\item Evidence that neither objective helps: SupCon loses $7.8$ accuracy
points --- the one primary comparison the bootstrap resolves --- and the GRL
$15.8$, with a $10\times$ larger seed spread and most of its damage in the
classifier head rather than in the representation.
\item A measured explanation: no domain attribute is linearly readable from
train features, replicated on five encoders and two rotations, and the
dataset's environment incidence makes this structural rather than accidental.
\item A statistical protocol for small LOEO test sets, under which the
backbone alone is worth $23.5$ accuracy points --- more than any objective we
tested.
\end{itemize}

\secref{sec:related_work} reviews related work,
\secref{sec:processing_architecture} and \secref{sec:model} the pipeline,
signals and splits, \secref{sec:learning_framework} the learning framework,
\secref{sec:results} the results, and \secref{sec:conclusions} concludes.
